\section{A unit native state scale: the same circuit with a different mass}

\begin{theorem}[Unit $\eta$ and unbounded $e$]
Use the unchanged physical native C1 source, its original event signs,
and the positive shapes and gaps $a,t,M,J$. Suppose $k,\zeta$ are
positive infinitesimals, $\eta$ is a positive bounded unit,
$r=1-a\eta$ is a positive bounded unit, and $e=k\sqrt{\eta/\zeta}$
is unbounded. Then the full original parent is tensor SOS at every
independent projective test and every real or complex state.
No positive lower bound on $a,t,M,J$ is imposed.
\end{theorem}

Write $H=r+\eta$ and $\Delta=1/e$. The seven native coquotients
have residues $2rH$ for $c_2,c_b,c_3,c_5,c_\alpha$, $2$ for
$c_4$, and $4r$ for $c_d$. All are positive units. In particular the
four retained packets of the contact--conductor proof remain positive
and the remainder has bounded coefficients. We reuse their exact
definitions, not asymptotic substitutes for the original source.

\paragraph{Outer tests.}
The exact formulas give
\[
 \frac md=\frac{e^2c_2}{M\eta c_d},\qquad
 \frac\beta d=\frac{e^2c_b}{t\eta c_d}-\frac12.
\]
Both ratios are unbounded. Apply the exact full-parent Alice reflection
(AS3)--(AS4), which exchanges $m\leftrightarrow c$ and
$\beta\leftrightarrow D$ while leaving $v$ unchanged. The already
proved two-strong-current identity then covers every test with $1/v$
bounded, for every Alice angle and every state. Its homogeneous form
includes both literal Alice endpoints, so no inverse-angle point is lost.

For $v$ infinitesimal and Alice infinity, use the critical-current
identity with $g_c=cD/(c+D)$, $\epsilon=d/g_c$ and
$\gamma=Ag_c/d^2$. Its actual sufficient conditions on this chart
are bounded $\epsilon$ and positive-unit $\gamma$: in its exact
four-root identity the correction $L_0$ has factors $c_u$ or $v$,
where $c_u=(1+u^2)^{-1/2}$ is infinitesimal. Thus $\sqrt\epsilon$
need only be bounded. The residue Gram is
$\operatorname{diag}(1,2,2,4\bar\gamma)$.
Use the exact relation $J+a=T$ with $T-1$ infinitesimal, and the
cancelled original coefficient formula
\[
 \frac Da=\frac{2c_5/c_3-J\eta c_d/c_3}{J}.
\]
Its numerator has residue $2/\bar H$. The unchanged native formulas give
\[
 \frac d a\sim\frac{4\eta}{H},\quad
 \frac A a\sim\frac{2\eta^2}{rH},\quad
 \frac D a\sim\frac2{JH},\qquad
 \epsilon\sim\frac{2\eta(JH+a)}H,\quad
 \gamma\sim\frac{H}{4r(JH+a)}.
\]
The identity
$\overline{JH+a}=1+\bar\eta(1-\bar a)^2$ shows that the two required
conditions hold even when $a$ or $J$ is infinitesimal. The calculation
uses the cancelled harmonic coefficient $g_c/a$, so no lower bound
on an opening is inserted. Homogeneous angle normalization includes
literal Alice infinity. It remains to treat bounded $u$ and
infinitesimal $v$.

\paragraph{State entry.}
Normalize the state by a coordinate of maximal absolute value and put
$D_u=1+u^2$ and $w=(W+uY)/D_u$. If $w$ is a unit, the exact five-root
current identity (UW2)--(UW3) applies with $Q=W+uY=D_uw$ a unit.
Its first four diagonal residues are $8rH$ and its fifth is
$8D_u$; both negative corrections have factors $\Delta$ or
$\Delta^2\eta$. Thus this chart is strict.

If $w$ is infinitesimal, the nonnegative pure-testing core is a direct
exit. On the negative-core branch $Z$ is a unit. Indeed its residue
at $v=0$ is $2\{(X+uZ)^2+(uW-Y)^2\}$; if $Z,w$ and both currents
had zero residue, every normalized state coordinate would have zero
residue. Normalize $Z=1$. The exact core Schur identity gives $w>0$
after the usual positive-angle sign gauges. All transformations here
belong solely to the testing field.

\paragraph{The ordinary-$u$ chart needs no contact division.}
Use the exact source change
\[
 \delta=\Delta R/K,\quad b^2=\eta K/R,\quad K=1+k^2,
 \quad Q=K+ak^2b^2,\quad T_n=K^2-t\delta^2k^4Qb^2.
\]
All source denominators are positive units, $\delta,k$ are
infinitesimal and $b^2\sim\eta$. Let
\[
 A_0=wv,\quad B_0=\delta v,\quad C_0=\delta b^2uw,
 \quad D_0=\delta^2k^3b^2,\qquad x=B_0-C_0.
\]
The five-term exact scalar belongs to
$(A_0,B_0,C_0)^2+(D_0^2)$ and both complete first normals belong
to $(A_0,B_0,C_0)$. Collecting these displayed factors before taking
residues is legitimate over the bounded coefficient ring. At
$\delta=k=w=v=0$, keeping bounded $u$ free, the scalar Gram is
\[
 8rH D_u^2A_0^2+
 8D_u\{x^2+u^2rH B_0^2\}+8rH D_uD_0^2.
 \tag{UE1}
\]
The hard block on $(N,\delta M_B)$ is
$\operatorname{diag}(16rH,16rH+8\eta^2/D_u)$, and the complete
first normals have residues $0$ and $16\eta x$. The latter is the
coefficient of $\delta M_B$ in the full polynomial, so the symmetric
Gram off-diagonal entry between $\delta M_B$ and $x$ is $8\eta$.
The coupled $(\delta M_B,x)$ block has determinant
\[
 (16rH+8\eta^2/D_u)\,8D_u-(8\eta)^2=128D_urH>0.
 \tag{UE2}
\]
If $u$ is a unit, the remaining $B_0$ pivot in (UE1) is strict,
as are the $A_0,D_0,N$ pivots. Every omitted coefficient has zero
residue. This proves the chart by a bounded strict Gram on these
six roots, without any mixed inverse. Equation (UE1) follows directly
from $a\eta+r=1$: its $(B_0,C_0)$ part before collection is
$8D_u\{(1+u^2rH)B_0^2-2B_0C_0+C_0^2\}$.

\paragraph{The infinitesimal-$u$ chart uses the old conductor.}
Now $\delta,k,u,w,v$ are infinitesimal, while $b,a,t$ are bounded.
Use the unchanged roots
\[
 y=A_0-kC_0,\quad z=uB_0-kC_0,\quad h=kC_0-D_0,
 \quad q_1=a_1a_2-a_3a_4,
\]
and the second pairing $q_2=a_1a_3-a_2a_4$ of the four separated
radicals. Set
\[
 T_4=\sum_{i=1}^4a_i^4-4\prod_{i=1}^4a_i
     =y^2+z^2+h^2-2yz+2q_1^2,
 \qquad \rho=(1-ab^2)(1+(1-a)b^2)-1.
\]
Give $(\delta,k,b,a,t,u,w,v)$ weights $(1,1,0,0,0,1,2,3)$.
The root $x$ has weight four and the other soft roots have weight
five. Substitution at the exact contact preserves this grading.
Writing $E_0$ for the complete scalar correction relative to the
old leading polynomial, the five-term factorization and the two
complete normal formulas give
\[
 \begin{aligned}
 E_0&=8\rho T_4-8(\rho+b^2)k^2x^2+O_{\rm wt}(12),\\
 L_1&=O_{\rm wt}(6),\qquad
 L_2=16b^2x+O_{\rm wt}(6).
 \end{aligned}\tag{UE3}
\]
These identities require only multiplication of the short source
polynomials. In particular the residue correction is the single
factor $\rho$ multiplying the same circuit, rather than a new
polynomial positivity problem.

The unlocalized monic contact--conductor lift applies unchanged.
After the terms explicitly displayed in (UE3) are removed, scalar
coefficient weights are at least two and normal coefficient weights
are at least one. Therefore they all have zero residue. There is no
division by $k,b,\delta,a,t$ or a contact difference. The leading
strict soft Gram is now
\[
                  8x^2+8rH T_4.
\]
The two-pairing circuit identity supplies a strict Gram on
$(y,z,h,q_1,q_2)$, and $rH$ is a positive unit. The complete hard
block is $\operatorname{diag}(16rH,16rH+8\eta^2)$; its only
nonzero first-normal residue is $16\eta x$. Its coupling to the
$8x^2$ pivot has determinant $128rH>0$. Retain the old positive
$g_4^2$ square. The resulting eight-root remainder Gram is strict.

Finally restore all four positive packets, the positive square just
retained, and the positive pure-coefficient multiplier. The state and
angle cases above exhaust the independent tests. Complex states split
into their real and imaginary quadratic evaluations. The result
closes the entire stated source domain, rather than only an optimizing
or matching testing chart.

\paragraph{Scope.}
This theorem adds the positive-unit $\eta,r$ slice with unbounded
$e$. It does not cover infinitesimal $r$, unbounded $\eta$, bounded
$e$, or the other native complement rows. Those domains must not be
marked complete by this result.
